\documentclass{article}
\usepackage{amsmath,amsthm,amssymb}

\newtheorem{definition}{Definition}[section]
\newtheorem{lemma}[definition]{Lemma}
\newtheorem{theorem}[definition]{Theorem}
\newtheorem{proposition}[definition]{Proposition}

\title{Definitions, Lemmas, and Proofs: A Small Mathematical Study}
\author{Sample Research Group}
\date{}

\begin{document}
\maketitle

\begin{abstract}
This paper illustrates the hierarchy of definitions, lemmas, theorems,
propositions, equations, and proofs. The results concern convergence of
real sequences and provide a test case for parsers that identify formal
environments and preserve structural relationships.
\end{abstract}

\section{Preliminaries}

Mathematical writing uses named environments to distinguish statements
from their justifications. A definition introduces terminology, a lemma
establishes an auxiliary result, and a theorem presents a result of
primary interest. A proof explains why a statement follows from
assumptions and previously established facts.

\begin{definition}
A sequence $(x_n)_{n\geq 1}$ of real numbers converges to
$L\in\mathbb{R}$ if, for every $\varepsilon>0$, there exists an integer
$N$ such that
\begin{equation}
n\geq N \implies |x_n-L|<\varepsilon.
\end{equation}
We write $x_n\to L$ as $n\to\infty$.
\end{definition}

\begin{definition}
A sequence $(x_n)$ is constant if there is a real number $c$ such that
$x_n=c$ for every positive integer $n$.
\end{definition}

\section{Elementary Lemmas}

\begin{lemma}
Every constant real sequence converges to its constant value.
\end{lemma}

\begin{proof}
Let $x_n=c$ for all $n$. Choose any $\varepsilon>0$ and set $N=1$.
For every $n\geq N$,
\[
|x_n-c|=|c-c|=0<\varepsilon.
\]
Thus the definition of convergence is satisfied, and $x_n\to c$.
\end{proof}

\begin{lemma}
If a sequence converges to $L$, then it is bounded.
\end{lemma}

\begin{proof}
Since $x_n\to L$, choose $N$ so that $n\geq N$ implies
$|x_n-L|<1$. By the triangle inequality,
\[
|x_n|\leq |x_n-L|+|L|<1+|L|
\]
for all $n\geq N$. The finite set of initial values
$x_1,\ldots,x_{N-1}$ has a finite maximum absolute value. Taking the
larger of that value and $1+|L|$ gives a bound for the whole sequence.
\end{proof}

\section{A Limit Rule}

\begin{theorem}
If $x_n\to a$ and $y_n\to b$, then $x_n+y_n\to a+b$.
\end{theorem}

\begin{proof}
Let $\varepsilon>0$. Convergence of $x_n$ gives an integer $N_1$ such
that
\[
n\geq N_1\implies |x_n-a|<\varepsilon/2.
\]
Convergence of $y_n$ gives an integer $N_2$ such that
\[
n\geq N_2\implies |y_n-b|<\varepsilon/2.
\]
Let $N=\max(N_1,N_2)$. For $n\geq N$, the triangle inequality yields
\begin{align}
|(x_n+y_n)-(a+b)|
&=|(x_n-a)+(y_n-b)|\\
&\leq |x_n-a|+|y_n-b|\\
&<\varepsilon.
\end{align}
Therefore, $x_n+y_n$ converges to $a+b$.
\end{proof}

\begin{proposition}
If $x_n\to a$ and $c$ is a fixed real number, then $cx_n\to ca$.
\end{proposition}

\begin{proof}
If $c=0$, the sequence $cx_n$ is identically zero and the result is
immediate. If $c\ne 0$, let $\varepsilon>0$. Choose $N$ so that
$n\geq N$ implies $|x_n-a|<\varepsilon/|c|$. Then
\[
|cx_n-ca|=|c||x_n-a|<\varepsilon.
\]
This proves the proposition.
\end{proof}

\section{Structural Dependencies}

The proof of the sum theorem uses the definition of convergence and the
triangle inequality. The proposition uses the definition again, while
the boundedness lemma is an auxiliary result that is not needed for the
sum theorem. A document-processing system should not assume that every
earlier statement is a dependency merely because it appears earlier.

LaTeX source order also differs from semantic dependency. A command such
as \verb|\ref{lem:bounded}| can link a statement to another environment,
while nested environments and equations create containment relationships.
A robust parser should preserve environment boundaries and extract
reference links separately from parent-child structure.

\section{Conclusion}

Formal mathematical documents contain useful structural boundaries.
Definitions, lemmas, theorems, and proofs can be represented as nodes,
but a tree alone does not capture every logical dependency. This makes
mathematical text a useful test domain for structured retrieval and
selective cache management.

\end{document}
